Im Comedy-Sketch ``Seven words you can't say on TV''\footnote{\url{https://youtu.be/5ssJtD08vCc}} erklärt George Carlin,
welche Wörter man im Fernsehen niemals verwenden kann\footnote{Es sind die
Wörter \texttt{shit}, \texttt{piss}, \texttt{fuck}, \texttt{cunt}, 
\texttt{cocksucker}, \texttt{motherfucker} und \texttt{tits}.}.
Da heutzutage Medieninhalte auch von Software erzeugt werden, wird es nötig,
Prüfprogramme zu haben, welche zuverlässig erkennen können, ob ein Programm
einen Text generieren wird, welcher verbotene Wörter enthält.
Ein Medienunternehmen verlangt daher nach einem Tool, mit dem man
beliebige Programme scannen kann.
Das Tool soll herausfinden, ob das gescannte Programm verbotene Wörter
generieren wird.
Das Medienunternehmen hat allerdings Mühe, einen Anbieter für ein
solches Tool zu finden.
Woran könnte das liegen?

\begin{loesung}
Sie \textit{FORBIDDEN} die Eigenschaft einer Sprache, dass eines der 
verbotenen Wörter als Teilkette in einem Wort der Sprache vorkommt.
Diese Eigenschaft ist nicht trivial, denn 
\[
\begin{aligned}
L_1&=\{\texttt{motherfucker}\} &&\text{hat die Eigenschaft \textit{FORBIDDEN},}
\\
L_2&=\{\}&&\text{hat die Eigenschaft nicht.}
\end{aligned}
\]
Die beiden Sprachen können von einem Programm erkannt werden.
Die Eigenschaft \textit{FORBIDDEN} ist daher nicht trivial.
Nach dem Satz von Rice ist $\textit{FORBIDDEN}_{\text{TM}}$ nicht
entscheidbar, aber genau nach einem Entscheider dafür hat das
Medienunternehmen verlangt.
\end{loesung}

\begin{bewertung}
Eigenschaft \textit{FORBIDDEN} ({\bf P}) 1 Punkt,
zwei Sprachen ({\bf L}) 2 Punkte,
Eigenschaft ist nicht trivial ({\bf T}) 1 Punkt,
Anwendung des Satzes von Rice ({\bf R}) 2 Punkte.
\end{bewertung}


